\section{Current Bottleneck and Trace Performance}
\label{sec:bottleneck}

\subsection{Measured Decode Latency (With Trace)}

Table~\ref{tab:timing} shows the per-component timing breakdown for a single decode step on the 1$\times$4 Wormhole mesh with trace enabled, averaged over 16 steady-state steps (steps 3--18, after 2 warmup and 1 first-trace-replay step).

\begin{table}[h]
\caption{Decode step timing (Granite-4.0-h-tiny, 1$\times$4 mesh, fabric + trace enabled).}
\label{tab:timing}
\begin{tabular}{lrr}
\toprule
Component & Time (ms) & \% of total \\
\midrule
Mamba2 (36 layers) & 52--58 & 54--61\% \\
MoE + Shared MLP (40 layers) & 62--68 & 65--71\% \\
Attention (4 layers) & 1.8--2.0 & $\sim$2\% \\
\midrule
Total (trace replay) & $\approx$95 & \\
Throughput & \multicolumn{2}{c}{10.5--10.6 tok/s} \\
\bottomrule
\end{tabular}
\end{table}

With trace enabled, the single \texttt{execute\_trace} dispatch replaces $\approx$500 individual PCIe round-trips. The step time drops from 130--145\,ms (no trace) to $\approx$95\,ms, a $\approx$33\% improvement. Mamba2 remains the dominant contributor because its per-layer recurrent ops are compute- and DRAM-bound; MoE all-gather collectives overlap with compute during replay.

\subsection{Bypass-Mode Off-by-One Discovery}

During development, an early inline capture strategy (triggering trace capture inside the generation loop at step $t$) produced incoherent responses across all prompts. Step-by-step token comparison showed that WITH TRACE, step $t$'s top token matched NO TRACE step $t-1$ — a consistent one-step lag. This was traced to Metal's bypass mode:

\begin{enumerate}
\item \texttt{begin\_trace\_capture()} calls \texttt{set\_bypass\_mode(true)} on all device CQs.
\item All ops between \texttt{begin} and \texttt{end} are recorded into a host-side buffer but \textbf{not dispatched} to the device.
\item KV cache writes (\texttt{paged\_update\_cache}) and SSM state updates (\texttt{ttnn.assign}) inside the capture window are recorded but never executed.
\item The device's state at the end of capture is identical to its state before capture began.
\item The captured step's \texttt{execute\_trace} replay then processes the \emph{current} embedding with the \emph{previous} device state, producing the same output as step $t-1$.
\end{enumerate}

This is not a bug in Metal; it is the intended behaviour that allows allocations to be fixed. The fix is to capture the trace with a dummy input \emph{outside} the generation loop (Algorithm~\ref{alg:trace}), so the capture step never corrupts the autoregressive state.

\subsection{Latency Decomposition}

We decompose the decode step latency into four components:
\begin{equation}
T_{\mathrm{step}} = T_{\mathrm{dispatch}} + T_{\mathrm{compute}} + T_{\mathrm{fabric}} + T_{\mathrm{DRAM}}
\label{eq:roofline}
\end{equation}
Without trace, $T_{\mathrm{dispatch}}$ from $\approx$500 PCIe round-trips dominates, accounting for the majority of the 130--145\,ms step time.

With trace, $T_{\mathrm{dispatch}}$ collapses to $\approx$1\,ms (one replay), reducing step time to $\approx$95\,ms. The remaining time is split across $T_{\mathrm{compute}}$, $T_{\mathrm{fabric}}$ (all-gather collectives), and $T_{\mathrm{DRAM}}$ (weight streaming); further profiling inside the trace replay would be required to separate these contributions.

\subsection{Path Forward}

\textbf{Fused SSM kernel.} The \texttt{ssm\_update} kernel (Algorithm~\ref{alg:kernel}) reduces Mamba per-layer DRAM passes from 5 to 3 and saves 4 dispatches per layer. With trace, the dispatch saving is already collapsed; the primary remaining benefit is eliminating one 768\,KB intermediate DRAM write per layer, reducing $T_{\mathrm{DRAM}}$.

\textbf{Fabric throughput.} The all-gather collectives for expert-parallel MoE are a significant contributor to the remaining 95\,ms step time. Profiling the per-all-gather latency inside trace replay and exploring ring-topology scheduling are the next steps to reduce this overhead.
